\chapter{量子理想气体}
现在我们已经准备好了一切必要的工具来分析理想气体的统计性质。理想气体指的是：粒子间没有相互作用的气体。在全同性的假设下，粒子自然地分为 Bose 统计与 Fermi 统计两类。我们将分别分析这两类气体的性质。
\section{理想 Fermi 气体}
理想 Fermi 气体的巨配分函数的对数可以写成：
\begin{linenomath}\begin{equation}
    \ln \varXi = \frac{1}{h^3}\int d^3p d^3 q D(p,q)\ln \left( 1 + e^{-\beta(\varepsilon(p) - \mu)} \right).
\end{equation}\end{linenomath}
其中 $\varepsilon(p)$ 是粒子的能量，$\mu$ 是化学势，$D(p,q)$ 是相空间的态密度。对于三维空间中的 Fermi 气体来说，这个态密度是
\begin{linenomath}\begin{equation}
    D(p,q) = 2s+1,
\end{equation}\end{linenomath}
因此
\begin{linenomath}\begin{equation}
    \ln \varXi = \frac{2s+1}{h^3}\int d^3p d^3 q \ln \left( 1 + e^{-\beta(\varepsilon(p) - \mu)} \right).
\end{equation}\end{linenomath}
于是各能级上占据的粒子数为
\begin{linenomath}\begin{equation}
    \langle n(p) \rangle = \frac{2s+1}{e^{\beta(\varepsilon(p) - \mu)} + 1}.
\end{equation}\end{linenomath}
其他物理量可以由前面所述的方式得出。
\subsection{零温度下的理想经典 Fermi 气体}
在零温度下， Fermi 气体的分布函数退化为一个阶梯函数：
\begin{linenomath}\begin{equation}
    \langle n(p) \rangle = \begin{cases}
        2s+1, & \varepsilon(p) < \mu\\
        0, & \varepsilon(p) > \mu
    \end{cases}.
\end{equation}\end{linenomath}
因此在零温度下，所有能量小于化学势 $\mu$ 的态都被占据，而所有能量大于化学势 $\mu$ 的态都没有被占据.
粒子数与化学势的关系这时候是：
\begin{linenomath}\begin{equation}
    N = \frac{2s+1}{h^3}\int_{\varepsilon(p) < \mu} d^3p d^3 q.
\end{equation}\end{linenomath}
当我们采用经典色散关系
\begin{linenomath}\begin{equation}
    \varepsilon(p) = \frac{p^2}{2m},
\end{equation}\end{linenomath}
我们可以得到化学势 $\mu$ 与粒子数 $N$ 的关系：
\begin{linenomath}\begin{equation}
    N = \frac{2s+1}{h^3}\int_{\varepsilon(p) < \mu} d^3p d^3 q = \frac{2s+1}{h^3}\int_{p^2 < 2m\mu} d^3p d^3 q = \frac{4\pi(2s+1)}{3h^3} (2m\mu)^{\frac{3}{2}}V.
\end{equation}\end{linenomath}
因此粒子数密度 $n:= \frac{N}{V}$ 便和化学势 $\mu$ 的关系为
\begin{linenomath}\begin{equation}
    \mu = \frac{1}{2m}\left( \frac{6\pi^2 n}{2s+1} \right)^{\frac{2}{3}}.
\end{equation}\end{linenomath}
这个化学势 $\mu$ 也被称为 Fermi 能量 $\varepsilon_F.$ 量 $(\frac{6\pi^2 n}{2s+1})^{\frac13}$ 具有波矢量纲，因此记作 $k_F,$ 被称为 Fermi 波矢。我们可以看到， Fermi 能量与 Fermi 波矢的关系为
\begin{linenomath}\begin{equation}
    \varepsilon_F = \frac{\hbar^2k_F^2}{2m}.
\end{equation}\end{linenomath}
这时， Fermi 气体的内能是
\begin{linenomath}\begin{equation}
    U = V\int_{0}^{p_F}\frac{p^2}{2m} 4\pi p^2 dp  = \frac{3}{5}N\varepsilon_F.
\end{equation}\end{linenomath}
压强是：
\begin{linenomath}
\begin{equation}
    p = -\left( \frac{\partial U}{\partial V} \right)_{N} = \frac{2}{5}\frac{N}{V}\varepsilon_F.
\end{equation}
\end{linenomath}
因此对于零温度下的理想 Fermi 气体，我们有
\begin{linenomath}\begin{equation}
    U = \frac32 pV.
\end{equation}\end{linenomath}
\subsection{低温下 Fermi 气体的 Sommerfeld 展开}
在 $T$ 很低时，Fermi 分布与阶梯函数产生了偏离，但若 $T$ 足够低，理想上我们便可以将这个偏离展成 $T$ 的级数。这种技术的系统展开方法首先是由 Sommerfeld 提出的，因此被称为 Sommerfeld 展开。

核心问题是计算如下的积分：
\begin{linenomath}\begin{equation}
    I(T) = \beta \int_{0}^{\infty} \frac{f(\varepsilon) (2s+1)}{e^{\beta(\varepsilon - \mu)} + 1} d\varepsilon.     
\end{equation}\end{linenomath}
其中 $f(\varepsilon)$ 是一个在 $\varepsilon = \mu$ 处连续的函数。我们设无量纲数 $z = \beta \varepsilon,$ 化学势也无量纲化为 $z_F = \beta \mu.$ 这样整个积分就化成：
\begin{linenomath}\begin{equation}
    I(T) = \int_{0}^{\infty} \frac{f(k_BTz) (2s+1)}{e^{z - z_F} + 1} dz.
\end{equation}\end{linenomath}
现在，将 $f(k_BTz)$ 在 $z<0$ 处定义为 $0,$ 这样我们就可以将积分的下限改为 $-\infty$：
\begin{linenomath}\begin{equation}
    I(T) = \int_{-\infty}^{\infty} \frac{f(k_BTz) (2s+1)}{e^{z - z_F} + 1} dz.
\end{equation}\end{linenomath}
这个量在 $T\to 0$ 时的极限是
\begin{linenomath}\begin{equation}
    I(0) = \int_{-\infty}^{z_F} f(k_BTz) (2s+1) dz = \int_{0}^{\mu} f(\varepsilon) (2s+1) d\varepsilon.
\end{equation}\end{linenomath}
定义函数
\begin{linenomath}\begin{equation}
    g(z) = \frac{1}{e^{z - z_F} + 1} - \tilde{\Theta}(z-z_F),
\end{equation}\end{linenomath}
其中 $\tilde{\Theta}(z-z_F)$ 定义如下：
\begin{linenomath}\begin{equation}
    \tilde{\Theta}(z-z_F) = \begin{cases}
        1, & z < z_F\\
        \frac12, & z = z_F\\
        0, & z > z_F
    \end{cases}.
\end{equation}\end{linenomath}
这样我们就可以将 $I(T)$ 写成
\begin{linenomath}\begin{equation}
    I(T) = I(0) + (2s+1)\int_{-\infty}^{\infty} f(k_BTz) g(z) dz.
\end{equation}\end{linenomath}
现在，我们作 Laplace 变换
\begin{linenomath}\begin{equation}
    \tilde{f}(s) = \int_{0}^{\infty} e^{-sz} f(k_BTz) dz.
\end{equation}\end{linenomath}
以及
\begin{linenomath}\begin{equation}
    \tilde{g}(s) = \int_{-\infty}^{\infty} g(z) e^{sz} dz.
\end{equation}\end{linenomath}
于是可以有
\begin{linenomath}
        \begin{align}
            \int_{-i\infty}^{i\infty} \tilde{f}(s) \tilde{g}(s) ds &= \int_{-\infty}^{\infty} \left( \int_{0}^{\infty} e^{-sz} f(k_BTz) dz \right) \left( \int_{-\infty}^{\infty} g(z') e^{sz'} dz' \right) ds\\
            &= \int_{0}^{\infty} \int_{-\infty}^{\infty} \int_{-i\infty}^{i\infty} f(k_BTz) g(z') e^{s(z'-z)} dz dz' ds\\
            &= \int_{0}^{\infty} \int_{-\infty}^{\infty} f(k_BTz) g(z') \left( \int_{-i\infty}^{i\infty} e^{s(z'-z)} ds \right) dz dz'\\
            &= 2\pi  \int_{0}^{\infty} \int_{-\infty}^{\infty} f(k_BTz) g(z') \delta(z'-z) dz dz'\\
            &= 2\pi  \int_{0}^{\infty} f(k_BTz) g(z) dz.
        \end{align}
\end{linenomath}
下面我们求出 $\tilde{g}(s)$ 的表达式：
\begin{linenomath}
        \begin{align}
            \tilde{g}(s) &= \int_{-\infty}^{\infty} g(z) e^{sz} dz\\
            &= \int_{-\infty}^{\infty} \left( \frac{1}{e^{z - z_F} + 1} - \tilde{\Theta}(z-z_F) \right) e^{sz} dz\\
            &= e^{z_Fs}\int_{-\infty}^{\infty} \left( \frac{1}{e^{x} + 1} - \tilde{\Theta}(x) \right) e^{sx} dx\\
            &= e^{z_Fs}\int_{-\infty}^{\infty} \frac{e^{sx}}{e^x + 1} dx - \frac1s e^{z_Fs}.
        \end{align}
\end{linenomath}
我们可以通过留数定理求出 $\int_{-\infty}^{\infty} \frac{e^{sx}}{e^x + 1} dx$ 的值，现在直接给出结果：
\begin{linenomath}\begin{equation}
    \int_{-\infty}^{\infty} \frac{e^{sx}}{e^x + 1} dx = \frac{\pi}{\sin \pi s}.
\end{equation}\end{linenomath}
因此
\begin{linenomath}\begin{equation}
    \tilde{g}(s) = \frac{e^{z_Fs}}{s}\left( \frac{\pi s}{\sin \pi s} - 1 \right).
\end{equation}\end{linenomath}




